% Source: ACM_Templates.pdf pp. 40-66. daily-problems/10-13/P3384 and user supplied segment tree; https://cp-algorithms.com/data_structures/segment_tree.html
\subsection{线段树：区间加与区间和}
\label{ds:seg-lazy}
\AlgorithmInfo
    {支持整段加法，作为树剖区间维护的基础。}
    {建树 $O(n)$；修改、查询 $O(\log n)$；空间 $O(n)$。}
    {$n\ge1$；闭区间、1 下标；和、标记及长度乘增量均不溢出。}
\subsubsection*{原理}
懒标记 \InlineCode{tag} 表示已计入当前节点、还未传给儿子的统一增量。长度为 $len$ 的区间加 $v$ 后，$sum\gets sum+len\cdot v$、$tag\gets tag+v$。完整覆盖时只更新当前节点；要进入儿子前，先用 \InlineCode{pushdown} 下传；修改儿子后再 \InlineCode{pushup}。

\InlineCode{eval} 将一次操作完整作用到节点，\InlineCode{pushdown} 复用它。查询可能下传标记，但不改变数组代表的值。

\subsubsection*{使用教程}
构造时指定 $n$，把初值放进 \InlineCode{w[1..n]}，再调用 \InlineCode{build(1,1,n)}。\InlineCode{modify(1,L,R,val)} 给整段加 $val$；\InlineCode{query(1,L,R).sum} 取区间和。实现沿用树剖代码的接口，本节不取模，需要取模时可用下一节的仿射模板。
\begin{lstlisting}[style=mtycpp,numbers=none]
SegTree seg(5);
seg.w={0,1,2,3,4,5};
seg.build(1,1,5);
seg.modify(1,2,5,3);
long long ans=seg.query(1,1,3).sum;
\end{lstlisting}
\needspace{7\baselineskip}
\subsubsection*{模板代码}
\CodeFile{code/data-structures/seg_lazy.cpp}
